\begin{table}[t]
\centering
\small
\begin{adjustbox}{max width=\linewidth}
\begin{tabular}{lrrrrr}
\toprule
$\sigma$ & $M$ & $K$ & RERANK gain & HOLDOUT gain & Optimism \\
\midrule
0.00025 & 10 & 5 & +1.35 & +0.07 & +1.28 \\
 & 25 & 5 & +1.34 & -0.01 & +1.35 \\
 & 50 & 10 & +1.70 & -0.05 & +1.75 \\
 & 100 & 20 & +2.01 & -0.14 & +2.15 \\
 & 125 & 25 & +2.09 & -0.17 & +2.26 \\
\midrule
0.0005 & 10 & 5 & +1.60 & -0.03 & +1.63 \\
 & 25 & 5 & +1.65 & +0.00 & +1.64 \\
 & 50 & 10 & +2.15 & +0.01 & +2.14 \\
 & 100 & 20 & +2.61 & -0.00 & +2.61 \\
 & 125 & 25 & +2.74 & -0.01 & +2.75 \\
\midrule
0.001 & 10 & 5 & +1.89 & -0.07 & +1.96 \\
 & 25 & 5 & +2.02 & -0.00 & +2.03 \\
 & 50 & 10 & +2.69 & +0.02 & +2.67 \\
 & 100 & 20 & +3.25 & -0.02 & +3.27 \\
 & 125 & 25 & +3.41 & -0.01 & +3.42 \\
\midrule
0.002 & 10 & 5 & +3.33 & +0.84 & +2.50 \\
 & 25 & 5 & +3.74 & +1.38 & +2.35 \\
 & 50 & 10 & +4.74 & +1.61 & +3.13 \\
 & 100 & 20 & +5.58 & +1.78 & +3.80 \\
 & 125 & 25 & +5.84 & +1.82 & +4.02 \\
\bottomrule
\end{tabular}
\end{adjustbox}
\caption{POST HOC (appendix). Search budget within each radius: random subsets of the 125 audit candidates, $K=\max(5,\mathrm{round}(0.2M))$, identical splits across $\sigma$ (2,000 trials).}
\label{tab:radius-budget}
\end{table}
